\documentclass[11pt]{article}
\usepackage[T1]{fontenc}
\usepackage[utf8]{inputenc}
\usepackage{amsmath,amssymb}
\usepackage[margin=1in]{geometry}
\usepackage{hyperref}
\hypersetup{colorlinks=true,urlcolor=blue,linkcolor=blue}
\setlength{\emergencystretch}{4em}
\title{P03.3: Full Curve Envelopes and Actual Projective Factors}
\author{Fresh mathematical draft}
\date{4 October 2026}
\begin{document}
\maketitle
This is a fresh, unverified mathematical draft addressing P03.3 only. It proves a sufficient geometric sector and an analytic branch counterexample. It does not claim the original sharp classification. In particular, no unpublished book assertion in the input is used as a theorem.

\section{Precise result and conventions}

Let K be a complete non-Archimedean field with nontrivial valuation. All spaces below are good, separated and strictly analytic. Put

E\_X(f) = \{h : X -> R : h <= f, and h is locally Fubini–Study or semipositive real PL\}, P\_X(f) = sup E\_X(f).

PL and LFS have their published meanings in Gubler–Jell–Rabinoff, arXiv:2111.05741v3 (7 March 2025), Definitions 5.1 and 7.8. In particular, an LFS function is an R-PL function whose residue line bundle is semiample; it is semipositive. Thus this exact union is the entire semipositive real-PL family. We do not replace it by invariant functions, graph functions, an equicontinuous family, or the LFS-only family. Real constants belong to the family: choose alpha in K with 0 < |alpha| < 1 and express a constant as a real multiple of log|alpha|.

\textbf{Theorem A (affinoid curves).} If X is a compact strictly affinoid rig-smooth curve of pure dimension one, then, for every finite continuous f on X, P\_X(f) is finite and continuous and P\_X(f)\textasciicircum{}* = P\_X(f). There is a sequence of GLOBAL semipositive R-PL minorants b\_n <= f with

\[
0\le P_X(f)-b_n\le 2/n\quad\text{everywhere on }X.
\]

No genus, tree, rationality, or acyclicity condition is imposed. X may be disconnected; connected normal curves over a perfect K, in particular algebraically closed K, are included. For EVERY analytic extension F/K,

\[
P_{X_F}(f\circ\pi_{F/K})=P_X(f)\circ\pi_{F/K}.
\]

The upstairs envelope here also uses its entire competitor union. All these envelopes and pullbacks are continuous. Their classical pshness is a separate consequence of uniform approximation, not the reason continuity holds. No CLD-v2 classification is asserted.

\textbf{Theorem B (actual projective geometry).} Let X be as in Theorem A. Let either Y = P\_X(V) for an analytic locally free sheaf V of positive rank, or Y = X x\_K Q\textasciicircum{}an for a smooth geometrically connected proper projective K-variety Q. Let p:Y -> X be the actual analytic projection. For a finite continuous obstacle f on Y define

\[
g(x)=\min_{y\in p^{-1}(x)} f(y).
\]

Then g is finite continuous, and

\[
P_Y(f)=P_X(g)\circ p.
\]

The right side is continuous and equals its usc regularization. This conclusion uses every competitor on Y and every finite continuous obstacle, including noninvariant obstacles. It also commutes with every analytic scalar extension. More generally the displayed envelope identity is valid for these projections over any compact good separated strict base; continuity then follows whenever the base has the asserted envelope-continuity property.

\textbf{Additional normalization sector.} Over a perfect nontrivially valued K, let X be a reduced strict affinoid curve whose normalization nu:Xtilde -> X is pointwise bijective. Then P\_X(f) is continuous for every finite continuous f. This conclusion is on the ground space. It does not assert that pointwise-bijective normalization persists under extension, or that this extra sector has the scalar-envelope identity of Theorem A.

\section{Checked public foundations and their precise uses}

Write GJR for arXiv:2111.05741v3; GRa for arXiv:2506.13548v2 (17 September 2025); GRb for arXiv:2503.05513v1 (7 March 2025); Thu for Thuillier's 2005 thesis, HAL tel-00010990v1; Ber for Berkovich, IHES 78 (1993).

\begin{enumerate}
\item GJR Definition 5.1 and Remark 5.2 define real PL by G-local finite real linear combinations of logarithms of analytic units. PL membership is a G-sheaf. Semipositivity is instead an ordinary analytic germ condition, GJR Definition 7.2 and Remark 7.3(3); it cannot merely be glued from positivity on a G-cover.
\item GJR 15.1–15.4 and Proposition 15.5 construct a vertical Cartier metric from a legally subdivided integral affine skeleton function. In the dense-value-group case subdivision vertices must lie in Sigma(Gamma), not at arbitrary real points. Proposition 15.5(2) says an R-PL function at a relative-boundary point of a rig-smooth curve is harmonic there; its affine reduction component imposes no nef degree sign.
\item In the proof of GJR Proposition 15.7, degree on an interior proper reduction component is the weighted sum of outgoing slopes. GJR Lemma 7.4 identifies nef formal line bundles and semipositive model functions. Genus affects the residue Picard class but not the nef degree condition on a curve.
\item GJR Proposition 7.6 gives analytic pullback, all analytic scalar extensions, and proper-surjective descent of semipositivity once R-PL membership is known. Corollaries 5.12 and 5.13 and Proposition 7.7 give real-PL and semipositive real-PL descent along finite Galois scalar extensions. Proposition 5.11 proves the PL descent by a finite-flat formal norm. We use these results with coefficient group R, which is divisible, and only finite extensions.
\item Thu Theorem 2.3.8 applies to a strict rig-smooth affinoid of pure dimension one. Its proof first embeds it, after finite separable extension, as an affinoid domain of a smooth projective curve. Thu Proposition 3.2.4 represents a Thuillier-lisse function on a compact analytic domain by a finite affine skeleton function after finite separable extension. Thu Proposition 3.2.7, its proof, Lemma 3.2.8, and Corollary 3.2.9 give uniform density on a compact smooth boundaryless curve. We apply those density statements on the ambient projective curve, not on a boundary point of the given affinoid. Thuillier-lisse membership is not itself our proof of real-PL membership.
\item GRa Theorem 4.3 says that on a good strict space over a nontrivially valued field an R-PL function is semipositive iff classically psh. Definition 4.1 quantifies over all analytic extensions and curve morphisms. Theorem 6.1 is the local maximum principle at interior points. Lemma 6.2 gives nonempty relative boundary for each positive-dimensional affinoid component. In 2.11(5'), extended boundary is contained over base boundary. Theorem 6.5 and its proof give stability under a pointwise limit that is already usc.
\item GRa Lemma 3.9(1,4,6) gives rig-smooth normalization over a perfect field, quotient topology for normalization, and R-PL descent through normalization. Its proof of (6) uses local constancy near a non-normal rig-point, so it does not require analytic invertible functions to generate every harmonic germ at a positive-genus vertex.
\item Ber Corollary 3.7.4 says a smooth morphism is open. The projective projections in Theorem B are smooth in this relative, boundaryless sense. A scalar projection is not declared open. Ber Theorem 7.1.1 gives algebraic/analytic cohomology comparison with compact support for a compactifiable scheme and an abelian torsion sheaf. Applied in degree zero to a proper connected scheme and the constant Z/2 sheaf, it proves connectedness of its analytification: a nontrivial analytic clopen partition would give an extra section. Properness makes compact-support degree-zero cohomology equal to ordinary degree-zero cohomology here.
\end{enumerate}

The complete literal selected pages, statement/proof selections, and checked applications are retained in source-extracts/source-provenance.json and the accompanying text files. GJR PDF page 78 and Thu PDF page 85 (printed page 71) were also inspected as rendered PDF images. The crucial GJR restriction is Gamma, not its unrestricted real closure; the crucial Thu density space is A\textasciicircum{}0, not a positivity cone. GRb's relevant statement in the supplied v1 is Proposition 5.1, although GRa's cross-reference calls its earlier numbering Proposition 3.1; we do not conflate these versions. Hashes only authenticate the supplied bytes.

\section{Actual real-PL realization on a finite skeleton}

First suppose Z is a compact rig-smooth curve with a strictly semistable formal model. Its skeleton Sigma is a finite weighted metric graph, and tau:Z -> Sigma is the deformation retraction. An edge e has length in Gamma = -log|K\textasciicircum{}*| and weight w(e). Boundary vertices B correspond to nonproper special-fiber components. For an affinoid, B is its relative boundary, and every connected component meets B.

A subdivision will be called legal here if every new cut is at a Gamma-coordinate on its edge. GJR 15.3 then gives an admissible formal subdivision model. The strictly semistable models in GJR 15.2 have smooth irreducible special-fiber components, hence no self-node loop edges. Cycles consisting of several edges are fully allowed.

\textbf{Realization lemma.} If A is continuous on Sigma, affine on each edge of a finite legal subdivision, with arbitrary REAL slopes and values, then H=A o tau is a GLOBAL R-PL function on Z. It is semipositive precisely when

\[
\sum_{e\text{ out of }v}w(e)\,\partial_e A\ge0
\]

at every interior vertex v of that subdivision. No inequality is required at B. It is harmonic away from vertices with nonzero slope sum (and is harmonic at relative-boundary vertices).

Here is the local formal construction, including the genus issue. Around a vertex v choose a small closed star with all arm lengths in Gamma. Its inverse image is a compact strict analytic neighborhood of v. After the legal formal subdivision, its special fiber consists of the central component C\_v and the arm components. For each arm i there is an integral function T\_i on this star which is distance from v on arm i and zero on every other arm.

On the formal node joining C\_v to arm i, use the semistable chart XY=gamma\_i, with -log|gamma\_i| equal to that arm's length and t=-log|X| measured from C\_v. Define the vertical Cartier divisor D\_i by the following local equations: X on the node chart, 1 on the central side away from that node, and gamma\_i on the arm side away from the node. The ratios are units: X is a unit on the central overlap; on the arm overlap X/gamma\_i=1/Y is a unit. On all other node charts the equation is 1. These local Cartier data also agree on etale overlaps, using the specified central/arm branches. This is the construction of GJR 15.4 in these local coordinates; for nonrational nodes the same Cartier data are taken on the etale chart and descend, and the closed-point residue degree is the edge weight.

The line bundle O(D\_i), trivialized on the generic fiber by its canonical section 1, has

\[
-\log\|1\|_{\mathcal O(D_i)}=T_i\circ\tau.
\]

This follows on each of the above charts from -log|X|, 0, and -log|gamma\_i|, respectively. Thus it is an actual PL function, G-locally the logarithm of an analytic unit with the appropriate sign. Its restriction to C\_v has divisor class O(p\_i), where p\_i is the attaching closed point; its degree is [kappa(p\_i):ktilde]=w(e\_i). This is why no PRINCIPAL divisor on C\_v is required. In positive genus O(p\_i) generally is not principal, and that causes no difficulty for a formal metric on the analytically trivial line bundle.

Now on this star

\[
A\circ\tau=A(v)+\sum_i s_i(T_i\circ\tau)
\]

for its real outgoing slopes s\_i. The constant is R-PL, so this is an actual finite real linear combination of PL metrics. On open edge neighborhoods an annulus coordinate gives the expression c+s(-log|T|). On every omitted residue disc the function is constant. These expressions cover Z and agree as functions, hence prove GLOBAL R-PL membership by Definition 5.1 and the sheaf property. They do not assert that all global real graph functions share one globally commensurable integral slope vector. There is no such requirement in real-PL membership.

At an interior vertex the residue real line bundle has degree sum\_i w(e\_i)s\_i. On the proper smooth reduction curve nef means degree >=0, and numerically trivial means degree zero. This proves the positivity and harmonicity assertions at ordinary germs, independently of G-cover gluing. At relative boundary the reduction component is nonproper and the nef curve test is vacuous, exactly as GJR Proposition 15.5(2). Off the finite vertices the formulas are harmonic. Points of types 1 and 4, including points not visible on the finite skeleton, lie in constant fibers locally; they introduce no untested sign or membership condition.

The construction permits nonprincipal degree-zero residue classes at genus vertices. It therefore does not confuse harmonic functions with ordinary-local logarithms of analytic units; GJR Remark 15.10 and Thu Theorem 2.3.21 explicitly distinguish them.

\section{The finite graph obstacle problem and full analytic comparison}

Let G be real affine on every edge of a finite legal subdivision of Sigma. For vertices V and a vector a in R\textasciicircum{}V, put

\[
D a(v)=\sum_{vw}\frac{w(vw)}{\ell(vw)}\bigl(a(w)-a(v)\bigr).
\]

Take all vectors satisfying a(v)<=G(v) and D a(v)>=0 for v outside B. Let m=min\_V G. The constant vector m is feasible. The coordinatewise maximum of two feasible vectors is feasible: at v select the vector supplying its maximum value, and each neighboring maximum is at least that selected neighbor value. Thus its D value cannot decrease at v.

Intersect the feasible set with [m,max\_V G]\textasciicircum{}V. This is compact and nonempty. Maximize each coordinate separately, and take the coordinatewise maximum of those finitely many maximizing vectors. The result u is the greatest feasible vector. It is greatest among ALL feasible vectors, since clipping another vector from below at m preserves feasibility.

At a boundary vertex, u(v)=G(v): increasing that coordinate would only help inequalities at adjacent interior vertices. At an interior vertex with u(v)<G(v), D u(v)=0: otherwise one could increase u(v) slightly, retaining its own inequality and helping its neighbors. Consequently the affine interpolant U of u is below G, has nonnegative interior slope sums, and is harmonic off the finite contact vertex set

\[
S=\{v:u(v)=G(v)\}.
\]

S contains B. By the realization lemma H=U o tau is a GLOBAL semipositive real-PL minorant of g=G o tau. A selected graph function has now been shown to be an actual global competitor.

To prove that it bounds the FULL analytic family, take any h in E\_Z(g). It is continuous and classically psh by GJR Definition 7.8 and GRa Theorem 4.3. On S, h<=g=H. Off S, H is harmonic in the real-PL sense, so h-H is classically psh there. If max\_Z(h-H)=M>0, the maximum locus is nonempty and closed in compact Z and disjoint from S. Every point of that locus lies in Int(Z), since B is contained in S. The local maximum principle makes that locus open in Z. It therefore contains an entire connected component of Z, contradicting that every such component meets B. Thus h<=H on ALL of Z, not only on the skeleton. We have proved

\[
P_Z(g)=H.
\]

This includes all noninvariant analytic competitors, every omitted direction, and type-4 points. It is not an envelope computation restricted to graph competitors.

The same comparison holds after EVERY analytic extension E/K. Pull H back to Z\_E. It remains semipositive, and harmonic off pi\textasciicircum{}\{-1\}(S), by GJR Proposition 7.6. Boundary satisfies partial Z\_E subset pi\textasciicircum{}\{-1\}(B), by GRa 2.11(5'), and every affinoid component of Z\_E has nonempty boundary, by Lemma 6.2. On pi\textasciicircum{}\{-1\}(S), H o pi=g o pi. Repeat the maximum-locus argument for an arbitrary upstairs competitor. This proves the strong exact identity

\[
P_{Z_E}(g\circ\pi_{E/K})=H\circ\pi_{E/K}.
\]

It requires neither openness of pi, nor any statement about commuting usc regularization, nor constancy of an unproved scalar fiber regularization.

\section{Uniform density for arbitrary obstacles and legal cuts}

For a strict rig-smooth affinoid X/K, finite separable extension embeds X in a smooth projective curve C; this is the first step in the checked proof of Thu Theorem 2.3.8. For disconnected X apply it componentwise and take a common finite extension. The projective Berkovich space is compact Hausdorff. The affinoid is a closed subspace, so Tietze's extension theorem extends any finite continuous obstacle f to a continuous real function on C\textasciicircum{}an.

Apply Thu Corollary 3.2.9 on this smooth BOUNDARYLESS curve, with the compact set X. For every eta>0 obtain a Thuillier-lisse function uniformly within eta of f on X. By Proposition 3.2.4, after another finite separable extension its restriction to X is A o tau for a continuous finite piecewise affine A on a semistable skeleton. This step gives graph factorization, not yet our PL realization. It is valid for arbitrary continuous f, not only radial or invariant f. In particular its uniform error controls values on all omitted directions and type-4 points.

A may have breakpoints not legal for a formal subdivision. Retain the original model vertices. On every edge, move each of the finitely many interior breakpoints to a nearby LEGAL coordinate and retain the corresponding function values; use affine interpolation on the moved intervals. If edge lengths are fixed and successive breakpoints stay distinct and in their original order, the new interpolant converges uniformly to the old one as the displacements tend to zero. For example, on each interval it is A composed with a piecewise affine increasing homeomorphism tending uniformly to the identity. Uniform continuity of A proves the asserted error estimate.

If Gamma is dense, coordinates in Gamma supply all needed approximations, without an extension. If Gamma is discrete, write Gamma=delta Z and choose a sufficiently fine grid (delta/N)Z, with N prime to char(K) when that characteristic is positive. Adjoin an N-th root of an element of valuation delta. This is a finite separable extension, and its value group contains that grid. Its semistable base-change model admits the desired legal Gamma-cuts. Thus arbitrary real subdivisions are never silently claimed to be formal-model subdivisions. In the algebraically closed case Gamma is divisible and dense, so no discrete case arises.

Finally pass to a finite Galois extension L/K containing these finitely many separable extensions. Scalar base change preserves the semistable model; its retraction and skeleton maps have the functoriality in Thu Proposition 2.2.21 and the node-chart description used in GJR 15.2. The pulled graph function is still a finite legal affine graph function, possibly on a split or weighted refined skeleton. Because scalar projection is surjective, all uniform error estimates are unchanged.

Choosing the preliminary density and interpolation errors with sum <=epsilon proves the following usable statement:

For every epsilon>0 there are a finite Galois L/K, a semistable model of X\_L, and a finite legal affine graph obstacle g=G o tau on X\_L with ||g-f o pi\_L|| <= epsilon.

No restriction has been placed on f. The realization lemma gives actual global real-PL membership for this g (it need not be semipositive); Section 4 supplies its actual semipositive graph envelope H=P\_\{X\_L\}(g).

\section{Finite Galois averaging, continuity, and all scalar extensions}

For sigma in Gal(L/K) put g\_sigma=g o sigma and H\_sigma=H o sigma. Each H\_sigma is a global semipositive R-PL function; all g\_sigma are within epsilon of f o pi\_L. Section 4's comparison applies to each translated pair and, importantly, after any further scalar extension.

Set A=(1/|Gal(L/K)|) sum\_sigma H\_sigma. It is invariant, continuous, globally semipositive R-PL, and A-epsilon <= f o pi\_L. By GJR Proposition 7.7 it descends to a GLOBAL semipositive R-PL function b\_epsilon on X with

\[
b_\epsilon\circ\pi_{L/K}=A-\epsilon.
\]

This uses sums, scalar multiplication, and verified invariant descent. It never assumes a maximum of real-PL functions is real PL.

Take ANY ground competitor h<=f. Then h o pi\_L-epsilon <= g\_sigma for every sigma. It is still a semipositive real-PL function, so exact full-family comparison implies h o pi\_L<=H\_sigma+epsilon for every sigma. Average and descend the inequality. Since b\_epsilon is a minorant,

\[
b_\epsilon\le P_X(f)\le b_\epsilon+2\epsilon.
\]

Taking epsilon=1/n gives the sequence in Theorem A. It proves finiteness (constants also give a baseline) and continuity as a uniform limit on the ENTIRE compact analytic space. Only now may one conclude P\_X(f)\textasciicircum{}*=P\_X(f).

For an arbitrary analytic extension F/K, take a common analytic extension E of F and the finite L; one can choose a field factor of F tensor\_K L, since L/K is finite separable. Pull an arbitrary full upstairs competitor h\_F<=f o pi\_F to X\_E. For every sigma its pullback minus epsilon is bounded by the pulled g\_sigma. The scalar form of Section 4 gives

\[
h_F\circ\pi_{E/F}\le(A+\epsilon)\circ\pi_{E/L}.
\]

Since A-epsilon=b\_epsilon o pi\_L, surjectivity of X\_E -> X\_F yields

\[
h_F\le b_\epsilon\circ\pi_{F/K}+2\epsilon.
\]

Conversely b\_epsilon o pi\_F is an upstairs competitor. Consequently

\[
b_\epsilon\circ\pi_{F/K}\le P_{X_F}(f\circ\pi_{F/K})\le b_\epsilon\circ\pi_{F/K}+2\epsilon.
\]

The SAME ground sequence tends uniformly to both sides' asserted common limit, proving the scalar identity in Theorem A. This handles all analytic extensions, not merely finite, algebraic, or completed algebraic-closure extensions.

For clarity, positivity of the resulting continuous function also has the full quantifiers. Each b\_n is classically psh by GRa Theorem 4.3, so its pullback along any extended analytic test curve is subharmonic. Uniform convergence persists under pullback. On a strict affinoid test domain, harmonic boundary comparison for the limit follows by applying it to b\_n with an added arbitrarily small constant and then passing to the limit. Equivalently this is the checked proof of GRa Theorem 6.5. This verifies the curve and scalar tests separately from the continuity proof. It asserts no CLD-v2 result.

If K is perfect, a normal curve is rig-smooth by GRa Lemma 3.9(1), so the theorem includes the requested arbitrary normal affinoids over algebraically closed K. Positive-genus reduction vertices and cyclic finite skeleta have nowhere been excluded.

For the extra normalization sector, a finite bijective normalization of an affinoid is a homeomorphism, by compactness and Hausdorffness (or GRa Lemma 3.9(4)). Given a semipositive real-PL function htilde on Xtilde, transfer it to a continuous function h on X. GRa Lemma 3.9(6) makes h R-PL. GJR Proposition 7.6(1') makes h semipositive, since normalization is proper surjective. Conversely every base competitor pulls back. Thus the two FULL envelopes agree through this homeomorphism. Theorem A on Xtilde gives continuity on X. Uniformly approximating competitors also descend. The statement does not identify envelopes on arbitrary scalar extensions where normalization can cease to be bijective.

In that last sector the ground envelope's pullback is still continuous and passes all extended curve tests, by its descended uniform semipositive approximants. What is deliberately left unclaimed is equality with the envelope of the enlarged upstairs competitor family.

\section{Proper projective factors and projective bundles}

First, a finite continuous classically psh function on a connected compact boundaryless analytic space is constant. It attains a maximum; by GRa Theorem 6.1 the maximum locus is open, and by continuity it is closed. Connectedness forces it to be the whole space. Hence on a connected proper projective analytic space the full obstacle envelope is simply the constant min f. This observation concerns ordinary real functions, not metrics with a nontrivial background line bundle.

This statement also applies to connected singular proper projective spaces: no normality assumption entered the maximum-locus proof. Thus normality is not a necessary condition in the original unrestricted domain class. The discontinuous affinoid branch example below and the constant proper envelopes must both be respected by any proposed sharp classification.

For p:Y -> X in Theorem B, the analytic fiber over x is a proper connected space over H(x): it is P(V\_x)\textasciicircum{}an or Q\_\{H(x)\}\textasciicircum{}an. Algebraic geometric connectedness makes the scheme over H(x) connected. To check the analytic connectedness used here, apply Ber Theorem 7.1.1 in degree zero with the constant Z/2 sheaf to this proper scheme. Algebraic global sections of that constant sheaf are exactly its two constant values; a nontrivial analytic clopen partition would give a third section, contradicting comparison. A competitor h on Y restricts to a classically psh function on this fiber, because the fiber is an analytic morphism into Y\_\{H(x)\} and the definition permits the scalar extension H(x)/K. It is continuous and therefore constant there. These are actual Berkovich fibers, not a product of sets of K-points.

Consequently h=a o p for a function a on X. Properness makes p a closed surjection with compact fibers; since h is continuous and fiberwise constant, the quotient topology makes a continuous. For a projective bundle, ordinary analytic trivializations of V give analytic sections over an ordinary open cover of X. Restrict h to such a section: a is R-PL and semipositive there, by GJR Proposition 7.6. The ordinary local definitions give GLOBAL semipositive R-PL membership of a.

For a factor Q with no K-rational point, choose a closed point of the nonempty projective variety Q. Its residue is a finite extension L/K, so Q\_L has an L-point. The section x -> (x,q) over X\_L shows a o pi\_L is semipositive R-PL. GJR Corollary 5.12 and Proposition 7.6(2') descend real-PL membership and semipositivity of a. No rational point over K is required and no unproved higher-dimensional PL descent is invoked.

Conversely every base semipositive real-PL a pulls back as an admissible competitor on Y. Since h<=f iff a(x)<=min\_\{p\textasciicircum{}\{-1\}(x)\} f=g(x), this proves a bijection between the entire competitor families. Therefore P\_Y(f)=P\_X(g) o p, including for the full exact union.

It remains to establish continuity of g, rather than assume it. Properness and compact fibers give attainment. If g(x)>c, the closed set \{y:f(y)<=c\} has closed image under p, and its complement is \{x:g(x)>c\}; hence g is lower semicontinuous. If g(x)<c, some y in the fiber has f(y)<c. The open set \{y:f(y)<c\} has open image, because p is smooth and thus open by Ber Corollary 3.7.4. Its image is \{x:g(x)<c\}; hence g is upper semicontinuous. It is therefore finite continuous. Theorem A now proves Theorem B.

For any analytic extension F/K and any x' over x, the extended analytic fiber is Y\_x tensor\_\{H(x)\} H(x'). Scalar projection of this fiber onto Y\_x is surjective. The minimum of f pulled back to it therefore is exactly g(x). Thus the new base obstacle is g o pi\_F. The family identification holds over F as well, and Theorem A gives the asserted raw-envelope scalar compatibility and continuity. Boundary is handled through the base curve; the projective projection has no relative boundary. There is no extra slope sign imposed on the inherited base-boundary fibers.

\section{A fresh actual analytic branch counterexample}

Let K be any complete nontrivially valued field, choose alpha with 0<|alpha|<1, put A=-log|alpha|, and fix a>0. Take

\[
X=\mathcal M\bigl(K\langle x,y\rangle/(xy)\bigr).
\]

It is the union of the closed x-disc D\_x and the closed y-disc D\_y, meeting at their rigid origins o. Define the finite continuous obstacle by f=0 on D\_x and

\[
f(y)=-a\min\{|y|,1\}=-a|y|
\]

on D\_y. The formulas agree at o, and their restrictions are continuous on the two closed branches. The y-Gauss point eta\_y has f(eta\_y)=-a, while f(o)=0.

For ANY full competitor h on X, its restriction to D\_y is semipositive R-PL and classically psh. The singleton-Shilov maximum principle, GRa Theorem 6.4, implies h<=h(eta\_y)<=-a on D\_y. The constant -a is a global competitor, so P\_X(f)=-a on the entire y-disc, including o.

For each integer n>=1 define the actual analytic function

\[
q_n=\log\max\{|\alpha|^n,|x|\},\qquad h_n=\frac{a}{nA}q_n.
\]

q\_n is PL on the global rational cover \{|x|<=|alpha|\textasciicircum{}n\}, \{|x|>=|alpha|\textasciicircum{}n\}: on the first domain it is n log|alpha|; on the second it is log|x| and x is invertible. Thus h\_n is a GLOBAL real-PL function. GRa Proposition 4.2 makes q\_n classically psh, as the finite maximum of a constant and log|x|, allowing the latter to be -infinity where x=0. GRa Theorem 4.3 makes h\_n semipositive. It is -a on D\_y and <=0 on D\_x, hence is a full competitor below f.

For any point of D\_x other than o, |x|>0, and h\_n ->0. Since f=0 there, P\_X(f)=0 on D\_x minus o. Therefore the raw envelope is exactly

\[
P_X(f)=\begin{cases}0&\text{on }D_x\setminus\{o\},\\-a&\text{on }D_y.\end{cases}
\]

Its usc regularization changes only o to 0. Every other y-branch point still has value -a. The sequence of y-rigid points y=alpha\textasciicircum{}m tends to o, proving an actual discontinuity. No book counterexample or previous generated proof is used.

The projective-family identification in Section 7 holds also over this singular base, without invoking Theorem A. Thus X x\_K P\textasciicircum{}r with obstacle pulled back from f has the pulled-back full raw envelope, and its regularization is likewise discontinuous at the nodal projective fiber. The last regularization statement follows locally from p being open: for a continuous open surjection p and locally bounded s, (s o p)\textasciicircum{}*=s\textasciicircum{}* o p, directly by the neighborhood definition of limsup. Alternatively evaluate along a fixed analytic projective section and the two branch approaches. These are actual proper projective factors; they do not cure the branch obstruction.

\section{What remains open and what this argument does not establish}

The original P03.3 asks for sharp additional geometric hypotheses on general good strict analytic domains, with necessity, all finite continuous obstacles, and the full competitor family. Theorems A and B do not classify normal or geometrically unibranch higher-dimensional spaces; they treat curves and geometry whose global competitors are forced to descend along proper fibers.

The new curve argument closes the realization obligation for legal finite-skeleton real functions on strict rig-smooth affinoids, with genus and cycles allowed. It does not prove a corresponding realization or density-preserving semipositive globalization theorem on arbitrary higher-dimensional normal domains. It proves no general envelope descent along a positive-dimensional nonprojective projection, arbitrary finite cover, or normalization with several points in a fiber. The extra pointwise-bijective-normalization sector is not asserted stable under arbitrary scalar extension.

Normal curves over imperfect fields which are not rig-smooth are not covered by the finite separable semistable-reduction argument as stated. Purely inseparable/perfection descent for their exact competitor family is not supplied here. No claim that algebraic closedness is necessary is made: Theorem A expressly removes it in the rig-smooth range. Relative-boundary regularity for general noncompact domains, general branched-curve necessity, and higher-dimensional necessity including codimension-two pinching remain unclassified. The fresh nodal example proves a concrete obstruction, not a necessity theorem for all nonnormal spaces.

No step infers continuity from positivity, an invariant toric result, a modulus bound on a selected family, or finite graph computations. The key comparison bounds EVERY analytic competitor, and the error is uniform for EVERY continuous obstacle and after EVERY analytic extension. This establishes the displayed sufficient sectors, while preserving the original target as open.

\section{Artifacts and reproducibility}

Inputs are retained verbatim in input.json and input-original.json; task-original.json preserves the task object and assignment-original.txt preserves the assignment string. The entire argument is also embedded in the standalone P03.3-paper.tex. status.json has only the requested keys.

Rerun source retention with: python prepare\_sources.py

Rerun document checks and the two TeX builds with: python verify\_deliverable.py

Actual command output, environment versions, source extraction diagnostics and build logs are under checks/. Rendered source checks are checks/gjr-78.png and checks/thuillier-85.png. The generated PDF is P03.3-paper.pdf when both installed-environment builds succeed. No package installation, alternate model or worker, external message, private campaign material, or heavy computation is used. Successful execution is administrative evidence, not mathematical verification.
\end{document}
